\documentclass[11pt,a4paper]{article}

\usepackage[a4paper,top=2.5cm,bottom=2.5cm,left=2cm,right=2cm]{geometry}
\usepackage{amsmath,amssymb,amsfonts,amsthm,mathtools}
\usepackage{booktabs}
\usepackage{enumitem}
\usepackage{microtype}
\usepackage[T1]{fontenc}
\usepackage{lmodern}
\usepackage[colorlinks=true,linkcolor=blue,citecolor=blue,urlcolor=blue]{hyperref}

\hypersetup{
  pdftitle={Chern--Simons Phase Selection on L(6,1): Differential-Cohomological Result and Equivariant-Lift Obstruction},
  pdfauthor={Damyan Damyanov},
  pdfsubject={ITT Paper 1: differential Chern--Simons phase, cyclic character, and equivariant-lift obstruction}
}

\newtheorem{definition}{Definition}[section]
\newtheorem{lemma}[definition]{Lemma}
\newtheorem{proposition}[definition]{Proposition}
\newtheorem{theorem}[definition]{Theorem}
\newtheorem{corollary}[definition]{Corollary}
\newtheorem{postulate}[definition]{Postulate}
\newtheorem{remark}[definition]{Remark}

\title{\textbf{Chern--Simons Phase Selection on $L(6,1)$:}\\
\textbf{Differential-Cohomological Result and Equivariant-Lift Obstruction}}
\author{\textbf{Damyan Damyanov}\\
\small\textit{Independent Researcher, Bulgaria}\\
\small\href{mailto:damione1@gmail.com}{damione1@gmail.com}}
\date{Research Note v0.1 --- October 2026}

\begin{document}
\maketitle

\begin{abstract}
We isolate the level-selection step of Information-Topological Theory (ITT) and formulate it in a mathematically conservative way. Starting from the lens-space sector $L(6,1)=S^3/\mathbb Z_6$ and the convention-fixed flat representation $\rho_1$, we use the differential Chern--Simons refinement of the integral level-$k$ class $k c_2$ to derive the closed-manifold phase
\[
Z_k(L(6,1),\rho_1)=\exp\!\left(2\pi i k\,\mathrm{cs}(\rho_1)\right)=e^{-2\pi i k/6}.
\]
This phase canonically induces, relative to the selected generator of $\mathbb Z_6$, a cyclic $U(1)$ character. We then analyze whether the same phase can be identified with the action of that generator on a hypothetical one-dimensional equivariant fiber over the three-dimensional connection sector. Two independent obstructions are found. First, the standard Chern--Simons line-bundle construction is a transgression to even-dimensional manifolds; the three-dimensional closed-manifold theory supplies a phase rather than an ordinary prequantum line over the three-dimensional connection space. Second, whenever a one-dimensional equivariant linearization exists, its fiber character is ambiguous up to $\operatorname{Hom}(\mathbb Z_6,U(1))\cong\mu_6$ unless an additional normalization datum is specified. Thus the canonical identification of the fiber character with the lens-space Chern--Simons phase does not follow from the presently specified data.

We therefore replace the unresolved equivariant-line identification by the explicit ITT phase-triviality selection postulate $Z_k(L(6,1),\rho_1)=1$. Under this postulate and ordinary integer level quantization, the minimal positive level is derived as $k_{\min}=6$. The result is not presented as a universal theorem of ordinary $SU(2)$ Chern--Simons theory.
\end{abstract}

\section{Introduction and Logical Status}

The purpose of this paper is to isolate one narrow mathematical issue in Information-Topological Theory (ITT): how the lens-space sector $L(6,1)=S^3/\mathbb Z_6$ constrains the Chern--Simons level.

The foundational formulation distinguishes ordinary Chern--Simons quantization,
\begin{equation}
k\in\mathbb Z,
\end{equation}
from the ITT-specific selection of a minimal level. The latter is based on the fundamental flat sector
\begin{equation}
\rho_1:\pi_1(L(6,1))\cong\mathbb Z_6\longrightarrow SU(2),
\end{equation}
with
\begin{equation}
\rho_1(\gamma)=
\begin{pmatrix}
 e^{2\pi i/6}&0\\
 0&e^{-2\pi i/6}
\end{pmatrix}
\end{equation}
and convention-fixed Chern--Simons invariant
\begin{equation}
\mathrm{cs}(\rho_1)=-\frac16\pmod{\mathbb Z}.
\end{equation}

Earlier formulations attempted to identify
\begin{equation}
\exp\left(2\pi i k\,\mathrm{cs}(\rho_1)\right)
\end{equation}
directly with a stabilizer character on a prequantum line fiber over a three-dimensional connection space. The present paper shows that this identification is not a theorem of the currently specified data. The correct derived result is the Chern--Simons phase itself, together with a precisely stated obstruction to promoting it canonically to the desired fiber action.

\section{Geometric Input}

\subsection{The lens space and its deck group}

Let
\begin{equation}
\Gamma=\mathbb Z_6=\langle\zeta\mid\zeta^6=1\rangle,
\qquad
\omega=e^{2\pi i/6}.
\end{equation}
The free action on
\begin{equation}
S^3=\left\{(z_1,z_2)\in\mathbb C^2:\ |z_1|^2+|z_2|^2=1\right\}
\end{equation}
is
\begin{equation}
\delta_\zeta(z_1,z_2)=(\omega z_1,\omega z_2).
\end{equation}
Hence
\begin{equation}
M:=S^3/\Gamma=L(6,1),
\qquad
\pi_1(M)\cong\Gamma.
\end{equation}

\subsection{The selected flat representation}

Fix
\begin{equation}
\rho_1(\zeta)=
\begin{pmatrix}
\omega&0\\
0&\omega^{-1}
\end{pmatrix}\in SU(2).
\end{equation}
The associated principal bundle can be realized equivariantly on the universal cover as
\begin{equation}
\widetilde P=S^3\times SU(2)
\end{equation}
with left $\Gamma$-action
\begin{equation}
\Phi_\gamma(x,h)
=
\left(\delta_\gamma x,\rho_1(\gamma)^{-1}h\right)
\end{equation}
and the usual right $SU(2)$ action on the second factor. The quotient
\begin{equation}
\boxed{
P_{\rho_1}=\widetilde P/\Gamma
}
\end{equation}
is a principal $SU(2)$ bundle over $L(6,1)$.

The trivial connection on $\widetilde P$ is $\Gamma$-invariant and descends to a flat connection $A_{\rho_1}$ with holonomy $\rho_1$. The non-trivial information is therefore carried by the quotient/equivariant descent data, not by a non-trivial flat connection on simply connected $S^3$.

\section{\texorpdfstring{The $\Gamma$-Equivariant Connection Groupoid}{The Gamma-Equivariant Connection Groupoid}}

Let
\begin{equation}
\mathcal A(S^3)=\Omega^1(S^3,\mathfrak{su}(2)).
\end{equation}
The connections on $\widetilde P$ that descend to $P_{\rho_1}$ are exactly
\begin{equation}
\boxed{
\mathcal A_{\rho_1}^{\Gamma}
=
\left\{
A\in\Omega^1(S^3,\mathfrak{su}(2)):\
\delta_\gamma^*A
=
\operatorname{Ad}_{\rho_1(\gamma)^{-1}}A,
\quad \gamma\in\Gamma
\right\}.
}
\end{equation}

The corresponding equivariant gauge group is
\begin{equation}
\boxed{
\mathcal G_{\rho_1}^{\Gamma}
=
\left\{
 u:S^3\to SU(2):
 u(\delta_\gamma x)
=
\rho_1(\gamma)^{-1}u(x)\rho_1(\gamma),
\quad \gamma\in\Gamma
\right\}.
}
\end{equation}
With
\begin{equation}
A^u=u^{-1}Au+u^{-1}du,
\end{equation}
we define the connection groupoid
\begin{equation}
\boxed{
\mathfrak C_{\rho_1}
:=
\left[\mathcal A_{\rho_1}^{\Gamma}/\mathcal G_{\rho_1}^{\Gamma}\right].
}
\end{equation}

The distinguished flat object is the class of the trivial connection on the covering space,
\begin{equation}
[A_0],
\qquad
A_0=0,
\end{equation}
together with its non-trivial equivariant lift determined by $\rho_1$.

\begin{remark}
The groupoid above is not restricted to the discrete flat moduli space. Continuous auxiliary paths relevant to spectral-flow or transgression questions live in the full equivariant connection space. This avoids identifying a continuous interpolation with a path inside a discrete set of flat equivalence classes.
\end{remark}

\section{Differential Chern--Simons Data}

\subsection{Integral characteristic class}

At integer level $k$, the relevant characteristic class is
\begin{equation}
kc_2\in H^4(BSU(2);\mathbb Z).
\end{equation}
Its differential refinement is a degree-four Cheeger--Simons differential character, denoted schematically by
\begin{equation}
\widehat c_{2,k}\in\widehat H^4(BSU(2);\mathbb Z).
\end{equation}
Differential characters refine integral characteristic classes by retaining both curvature and secondary holonomy data\cite{CheegerSimons1985}.

Pulling the differential class back along the classifying data of $(P_{\rho_1},A_{\rho_1})$ gives
\begin{equation}
\widehat C_k(P_{\rho_1},A_{\rho_1})
\in
\widehat H^4(L(6,1);\mathbb Z).
\end{equation}

\subsection{Three-dimensional evaluation}

For a closed oriented three-manifold, a degree-four differential character evaluates on the fundamental three-cycle and produces a $U(1)$ phase. In the Chern--Simons case this evaluation is the exponentiated Chern--Simons invariant\cite{Freed1995,CheegerSimons1985}:
\begin{equation}
\boxed{
Z_k(M,A)
=
\operatorname{Hol}_{M}\left(\widehat C_k(P,A)\right)
=
\exp\left(2\pi i k\,\mathrm{cs}(A)\right).
}
\end{equation}

For the selected lens-space sector,
\begin{equation}
\mathrm{cs}(\rho_1)=-\frac16\pmod{\mathbb Z},
\end{equation}
so
\begin{equation}
\boxed{
Z_k(L(6,1),\rho_1)=e^{-2\pi i k/6}.
}
\end{equation}

\begin{proposition}[Differential-CS phase]
The phase
\begin{equation}
Z_k:=Z_k(L(6,1),\rho_1)
\end{equation}
is determined by the standard differential Chern--Simons refinement together with the fixed convention for $\mathrm{cs}(\rho_1)$. No ITT-specific equivariant-line assumption is needed for this result.
\end{proposition}

\begin{proof}
The differential refinement represents the integral class $k c_2$. Evaluation on the closed three-manifold gives the exponentiated Chern--Simons invariant. Substitution of the fixed lens-space value $\mathrm{cs}(\rho_1)=-1/6$ gives the stated phase.
\end{proof}

\section{The Induced Cyclic Character}

Because $k\in\mathbb Z$,
\begin{equation}
Z_k^6=e^{-2\pi i k}=1.
\end{equation}
Therefore, after fixing the generator $\zeta$, the phase defines a genuine cyclic character
\begin{equation}
\boxed{
\chi_k^{\rm CS}:\Gamma\to U(1),
\qquad
\chi_k^{\rm CS}(\zeta^n)=Z_k^n.
}
\end{equation}
Thus
\begin{equation}
\boxed{
\chi_k^{\rm CS}(\zeta)=e^{-2\pi i k/6}.
}
\end{equation}

This construction is elementary once $Z_k$ is known. It should not be conflated with a theorem identifying $\chi_k^{\rm CS}$ with the fiber character of an independently defined equivariant anomaly line.

\begin{remark}[Restriction to the selected sector]
The construction above is intentionally attached to the distinguished sector $\rho_1$. For general lens-space sectors $\rho_n$, the Chern--Simons invariant is quadratic in the holonomy label, whereas a group character is linear under powers of a fixed generator. Therefore no general identity of the form $\mathrm{cs}(\rho_n)=\chi(\zeta^n)$ is asserted.
\end{remark}

\section{Why the Canonical Equivariant-Line Identification Fails}

\subsection{Dimension of the standard Chern--Simons line construction}

The classical geometric Chern--Simons theory constructs a prequantum line bundle with connection over the moduli space of flat connections on a \emph{two-dimensional} manifold\cite{Freed1995}. More generally, the equivariant Chern--Simons bundle construction of Ferreiro Pérez is formulated for compact even-dimensional underlying manifolds and uses equivariant holonomy data\cite{Ferreiro2019}.

The three-dimensional closed-manifold theory instead assigns a $U(1)$ phase to the full three-dimensional field. Consequently, the expression
\begin{equation}
\mathcal L_k\longrightarrow
\mathcal A(L(6,1))
\end{equation}
as a standard three-dimensional prequantum line bundle is not a legitimate use of the ordinary CS line-bundle construction.

The appropriate three-dimensional datum may be represented by the Chern--Simons phase itself or, in higher-geometric language, by the corresponding invertible/anomaly theory. But this does not automatically produce the desired one-dimensional stabilizer fiber character.

\subsection{Equivariant linearization ambiguity}

Suppose nevertheless that a one-dimensional complex fiber
\begin{equation}
\mathcal L_{k,\rho_1}
\end{equation}
has been equipped with a genuine $\Gamma$-equivariant structure. The generator acts by
\begin{equation}
\widehat\Phi_\zeta
=
\chi_k^{\rm anom}(\zeta)\,\mathrm{id}.
\end{equation}

Let
\begin{equation}
\lambda:\Gamma\to U(1)
\end{equation}
be any character. Then
\begin{equation}
\boxed{
\widehat\Phi'_\gamma
=
\lambda(\gamma)\widehat\Phi_\gamma
}
\end{equation}
is again a genuine equivariant structure. Indeed,
\begin{align}
\widehat\Phi'_{\gamma_1}\widehat\Phi'_{\gamma_2}
&=
\lambda(\gamma_1)\lambda(\gamma_2)
\widehat\Phi_{\gamma_1}\widehat\Phi_{\gamma_2}
\\
&=
\lambda(\gamma_1\gamma_2)
\widehat\Phi_{\gamma_1\gamma_2}
\\
&=
\widehat\Phi'_{\gamma_1\gamma_2}.
\end{align}

For $\Gamma=\mathbb Z_6$,
\begin{equation}
\operatorname{Hom}(\Gamma,U(1))\cong\mu_6.
\end{equation}
Hence the fiber character is not uniquely determined by the underlying line and ordinary differential Chern--Simons class alone.

\begin{proposition}[Character ambiguity]
The data
\begin{equation}
\left(L(6,1),\rho_1,\widehat c_{2,k}\right)
\end{equation}
do not uniquely determine a $\mathbb Z_6$-equivariant linearization of an arbitrary one-dimensional fiber. Any proposed character can be modified by a factor in $\operatorname{Hom}(\mathbb Z_6,U(1))\cong\mu_6$ without changing the underlying non-equivariant fiber.
\end{proposition}

\begin{proof}
The construction above explicitly twists an equivariant structure by an arbitrary character $\lambda$. The resulting linearization obeys the same group law. Therefore an additional normalization datum is necessary to single out one character.
\end{proof}

\subsection{The mapping-torus route does not recover the lens-space phase}

One might try to obtain the desired generator action from an ordinary mapping torus. Consider the smooth isotopies
\begin{equation}
\delta_t(z_1,z_2)
=\left(e^{2\pi i t/6}z_1,e^{2\pi i t/6}z_2\right)
\end{equation}
and
\begin{equation}
h_t=
\begin{pmatrix}
 e^{2\pi i t/6}&0\\
 0&e^{-2\pi i t/6}
\end{pmatrix},
\end{equation}
with $0\le t\le1$. These give an isotopy of the lifted bundle automorphism from the identity to
\begin{equation}
\Phi_\zeta(x,u)
=
\left(\delta_\zeta x,\rho_1(\zeta)^{-1}u\right).
\end{equation}
Thus the associated ordinary mapping-torus construction is isotopic to the product construction on $S^3\times S^1$ and cannot reproduce the non-zero rational invariant $-1/6$ as an ordinary closed-four-dimensional second Chern number.

This does not contradict the lens-space Chern--Simons invariant. The latter is a three-dimensional secondary invariant of the quotient flat sector, whereas the mapping-torus characteristic number measures a different four-dimensional construction.

\begin{remark}
The distinction also rules out replacing the rational lens-space invariant by the assertion
\begin{equation}
\int_{L(6,1)\times S^1}c_2=-\frac16.
\end{equation}
For an honest closed four-manifold and an honest $SU(2)$ bundle, the ordinary characteristic number is integral.
\end{remark}

\section{D-01b.3: Precise No-Go Statement}

The preceding arguments establish the following restricted conclusion.

\begin{theorem}[No canonical three-dimensional fiber identification from the present data]
From the specified data
\begin{equation}
\left(L(6,1),\rho_1,\widehat c_{2,k}\right)
\end{equation}
one can derive the canonical Chern--Simons phase
\begin{equation}
Z_k=e^{-2\pi i k/6}
\end{equation}
and the induced cyclic character
\begin{equation}
\chi_k^{\rm CS}(\zeta)=e^{-2\pi i k/6}.
\end{equation}
However, those data do not by themselves canonically determine an identification
\begin{equation}
\mathcal L_{k,\rho_1}^{\Gamma}
\cong
\mathbb C_{\chi_k^{\rm CS}}
\end{equation}
for a hypothetical one-dimensional equivariant fiber over the three-dimensional connection sector. Such an identification requires additional structure or normalization beyond the currently specified differential Chern--Simons data.
\end{theorem}

\begin{proof}
The differential character determines the closed-manifold phase by standard Chern--Simons evaluation. The standard line-bundle construction is a transgression associated with even-dimensional base manifolds, so it does not by itself supply a canonical ordinary line over the three-dimensional connection space. Independently, any one-dimensional equivariant linearization is ambiguous under twisting by an arbitrary $U(1)$ character of $\Gamma$. Hence there is no canonical identification with the specific character $\chi_k^{\rm CS}$ without additional data.
\end{proof}

\section{Revised ITT Selection Principle}

The obstruction above does not remove the derived Chern--Simons phase. Instead, it changes the location of the ITT-specific input.

\begin{postulate}[ITT Fundamental-Sector Phase Triviality]
The fundamental ITT quotient sector is admissible at level $k$ only when
\begin{equation}
\boxed{
Z_k(L(6,1),\rho_1)=1.
}
\end{equation}
\end{postulate}

Using
\begin{equation}
Z_k=e^{-2\pi i k/6},
\end{equation}
we obtain
\begin{equation}
e^{-2\pi i k/6}=1
\quad\Longleftrightarrow\quad
\frac{k}{6}\in\mathbb Z.
\end{equation}
Thus
\begin{equation}
\boxed{k\in6\mathbb Z.}
\end{equation}

\begin{corollary}[Conditional minimal level]
Under ordinary integer Chern--Simons quantization and the ITT fundamental-sector phase-triviality postulate,
\begin{equation}
\boxed{k_{\min}=6.}
\end{equation}
\end{corollary}

\section{Interpretation and Limitations}

The mathematical chain established in this paper is
\begin{equation}
\boxed{
\mathrm{cs}(\rho_1)=-\frac16
\Longrightarrow
Z_k=e^{-2\pi i k/6}
\Longrightarrow
\chi_k^{\rm CS}(\zeta)=e^{-2\pi i k/6}.
}
\end{equation}
The next step,
\begin{equation}
\chi_k^{\rm CS}=1,
\end{equation}
is not forced by ordinary Chern--Simons theory. It is the explicit ITT phase-triviality selection postulate.

Accordingly, this paper does not claim a universal derivation of $k=6$ from $SU(2)$ Chern--Simons theory. It establishes the standard differential-CS phase, isolates the obstruction to the originally proposed canonical equivariant-fiber interpretation, and gives a mathematically unambiguous location for the ITT-specific input.

No argument involving $3\times2$, three fermion generations, Standard-Model anomaly counting, or phenomenological parameter fitting is used. The future generation problem remains independent.

\section{Conclusion}

The principal outcome is a clean separation between what is supplied by standard Chern--Simons geometry and what must be specified by ITT.

First,
\begin{equation}
\mathrm{cs}(\rho_1)=-\frac16
\end{equation}
is a convention-fixed rational invariant of the selected flat lens-space sector. Second, differential Chern--Simons theory derives
\begin{equation}
Z_k=e^{-2\pi i k/6}.
\end{equation}
Third, this phase defines a cyclic character relative to the chosen generator of $\mathbb Z_6$. Fourth, the desired identification of this character with a canonical stabilizer action on a one-dimensional three-dimensional prequantum fiber is not determined by the present data; the standard transgression dimension and equivariant character ambiguity provide independent obstructions.

The mathematically conservative ITT replacement is the explicit selection principle
\begin{equation}
Z_k(L(6,1),\rho_1)=1,
\end{equation}
which yields
\begin{equation}
\boxed{k_{\min}=6}
\end{equation}
as an ITT-specific consequence conditional on that postulate. This is the complete level-selection result established by the present paper.

\begin{thebibliography}{99}

\bibitem{CheegerSimons1985}
J. Cheeger and J. Simons,
\textit{Differential characters and geometric invariants},
in \textit{Geometry and Topology}, Lecture Notes in Mathematics \textbf{1167}, 50--80,
Springer (1985).

\bibitem{Freed1995}
D. S. Freed,
\textit{Classical Chern--Simons Theory, Part 1},
Adv. Math. \textbf{113}, 237--303 (1995).

\bibitem{Ferreiro2019}
R. Ferreiro Pérez,
\textit{Equivariant differential characters and Chern--Simons bundles},
arXiv:1907.00292 [math.DG] (2019).

\bibitem{Jeffrey1992}
L. C. Jeffrey,
\textit{Chern--Simons--Witten Invariants of Lens Spaces and Torus Bundles, and the Semiclassical Approximation},
Commun. Math. Phys. \textbf{147}, 563--604 (1992).

\bibitem{KirkKlassen1990}
P. Kirk and E. Klassen,
\textit{Chern--Simons Invariants of 3-Manifolds and Representations of Discrete Groups},
Math. Ann. \textbf{287}, 343--367 (1990).

\bibitem{Witten1989}
E. Witten,
\textit{Quantum Field Theory and the Jones Polynomial},
Commun. Math. Phys. \textbf{121}, 351--399 (1989).

\end{thebibliography}
\end{document}
